\documentclass[onecolumn]{article} %Este comando fabrica por defecto el tipo de documento, en este caso un artículo, en el formato de dos columnas
\usepackage[spanish]{babel}
\usepackage[left=2cm, right=2cm, top=2cm, bottom=2cm]{geometry} %Esto es para ponerle los márgenes, puedes hacerlo más grande o más chico
\usepackage{fancyhdr} %no borres esto, es para estilizar el encabezado y el pie de página
\usepackage{graphicx} %este es para adjuntar imágenes
\usepackage{float}
\usepackage{subfig} 
\usepackage{amssymb}
\usepackage{gensymb}
\usepackage{svg}
\usepackage[colorlinks=true, linkcolor=blue, urlcolor=blue, citecolor=blue]{hyperref}

% Configuración del encabezado solo en la primera página
\fancypagestyle{firstpage}{
	\pagestyle{fancy} 
	\fancyhead{}
	\fancyhead[L]{\textbf{Universidad de Antioquia}} % Fecha en la esquina izquierda
	\fancyhead[R]{\textbf{\today}} % Nombre de la institución en la derecha
	\fancyfoot[C]{\thepage} % Número de página centrado
}

\begin{document}
	
	\title{
		\textbf{Informe de lectura: Hábito 7 - Afilar la sierra}
		}
		
	
	\author{
		\centering
		\begin{minipage}{0.45\textwidth}
			\centering
			\textbf{David Alejandro Morón Acacio}\\
			Facultad de ingeniería\\
			\texttt{david.moron@udea.edu.co}\\
			\vspace{10pt}  
			\textbf{Luisa María Bohórquez Ardila}\\
			Facultad de ingeniería\\
			\texttt{luisam.bohorqueza@udea.edu.co}\\
		\end{minipage}%
		\hfill
		\begin{minipage}{0.45\textwidth}
			\centering
			\textbf{Iván René Acosta Lallemand}\\
			Facultad de ingeniería\\
			\texttt{ivan.acosta@udea.edu.co}\\
			\vspace{10pt}
			\textbf{Juan David Moreno Díaz}\\
			Facultad de ingeniería\\
			\texttt{juan.morenod@udea.edu.co}
		\end{minipage}
	}
	
	\date{} %este comando es necesario para borrar la fecha que maketitle pone por defecto
	
	\maketitle
	
	\thispagestyle{firstpage}
	
	\section{Introducción}
	
	El presente informe analiza el séptimo hábito del libro de Stephen Covey llamado "Los siete hábitos de la gente altamente efectiva". El objetivo principal es explorar cómo este hábito actúa como el motor de la mejora continua, permitiendo que un líder mantenga su efectividad a largo plazo sin sucumbir al agotamiento. Justificamos este estudio en la realidad actual de las organizaciones, donde la presión por resultados inmediatos a menudo hace que los gerentes descuiden su "capacidad de producción", es decir, a ellos mismos como su activo más valioso.
	
	\section{Resumen del capítulo IV}
	
	\section{Metodología o sinopsis según las secciones del capítulo}
	
	Para esta parte del informe se presenta la trama principal de cada sección, partes llamativas, eventos y características claves. 
	
	\subsection{Séptimo hábito: Afile la sierra - Principios de autorrenovación equilibrada}
	
	Inicialmente mente nos encontramos con una imagen que va desde la dependencia, pasa por la independencia, luego por la interdependicia y todos esto rodeado por un hábito llamo la interdependencia, entre la dependencia y la independencia hallamos la victoria privada que está conformada por los siguientes hábitos: sea proactivo, empiece con un fin en mente, 
	
	\subsection{Las cuatro dimensiones de la renovación}
	
	\subsubsection{La dimensión física}
	
	\subsubsection{La dimensión espiritual}
	
	\subsubsection{La dimensión mental}
	
	\subsubsection{La dimensión emocional/social}
	
	\subsection{Programando a los otros}
	
	\subsection{El equilibrio en la renovación}
	
	\subsection{La sinergia en la renovación}
	
	\subsection{La espiral asecendente}
	
	\subsection{Sugerencias prácticas}
	
	\subsection{Otra vez de dentro hacia afuera}
	
	\subsection{La vida intergeneracional}
	
	\subsection{Convirtiéndose en una persona de transición}
	
	\subsection{Una nota personal}
	
	\section{Análisis crítico y opinión personal}
	
	\section{Conclusión}
	
	\section{Trabajo de equipo en la actividad}
	
	\subsection{¿Nombre del líder cómo y por qué fue elegido?}
	
	La líder es Luisa María Bohórquez Ardila, a ella le apasiona el tema del desarrollo personal y se ofreció para ser la líder de la actividad. Las personas que estuvieron en la reunión el día que había que decidir el líder estuvieron de acuerdo por su entusiasmo por lo que ella quedó como líder.
	
	\subsection{¿Cómo se realizó la distribución o asignación de las actividades y responsabilidades?}
	
	Primero se realizó una lluvia de ideas y una consulta sobre las tareas que se querían realizar, teniendo en cuenta las preferencias de los participantes la líder desarrolló un plan de acción diario y semanal dónde se registraban los avances de las actividades y las fechas destinadas a las futuras tareas. También hubo al menos 4 controles por parte de la líder para verificar que todo estuviese en orden. Para tener mayor enfoque en las tareas la líder sugirió la conformación de pequeños equipos de trabajo para las tareas más complejas y los integrantes se unieron según sus preferencias, los grupos fueron para el informe y el guión de la obra de teatro (vídeo). Sin embargo la líder quería que todos se vieran involucrados en todas las actividades, entonces dentro del plan de acción todos los participantes debieron hacer su respectivo análisis de las páginas del libro que sirvieran tanto para el informe como para la obra de teatro, por lo que ambos microequipos tuvieron una base para trabajar. Finalmente, tanto el equipo del informe como el de la obra de teatro recibieron retroalimentación donde hubo correcciones y sugerencias por parte del resto de integrantes y así se logró una sinergia dentro del equipo.
	
	\subsection{¿Con cuál o cuáles de las habilidades gerenciales mencionadas a continuación, se relaciona principalmente el hábito y porque?}
	
	\subsubsection{Habilidades de inteligencia emocional}
	
	El hábito se relaciona con la autoconsciencia, la gestión de emociones y el manejo de conflictos de forma directa e intrínseca, pues a través de la autoconciencia y la gestión de las emociones se cumple el objetivo de manejar los conflictos efectivamente.
	
	\subsubsection{Habilidades de comunicación}
	
	Hay relación especial con la comunicación efectiva y la escucha activa, porque se tiene como una de las bases el factor emocional/social que nos lleva a prestar atención a las cosas que nos dicen y a hablar desde los hechos con las palabras adecuadas según el contexto.
	
	\subsubsection{Habilidades de relacionamiento}
	
	En este apartado podemos resaltar la estrecha conexión entre el sétimo hábito y el construir relaciones significativas, afilar la sierra habla desde uno mismo y hacia uno mismo, pero teniendo en cuenta las relaciones que nos nutren y que nos mantienen rodeados de personas que nos ayudan a crecer y preservar nuestra integridad, llenándonos de cosas buenas y de lazos que suman en lugar de restar.
	
	\subsubsection{Habilidades de logro de resultados}
	
	En este caso, es el alcance de metas el que tiene relación directa con nuestro hábito, afilar la sierra significa que debemos mantenernos a nosotros mismos enteros y dispuestos para seguir logrando nuestros objetivos, ya que si no cuidamos de nosotros mismos no podremos seguir avanzando.
	
	\subsubsection{Habilidades de desarrollo de otros}
	
	Hay una íntima conexión con la habilidad colaborativa y para trabajar en equipo de forma activa y eficaz, pues a través de la escucha y la comunicación se llega a ese aspecto de afilar la sierra que se nutre del buen compañerismo y de tenerse en cuenta unos a otros para resolver conflictos.
	
	\subsubsection{Habilidades de gestión}
	
	Aquí se toman dos habilidades principales: la toma de decisiones y el manejo del tiempo, pues ambas se benefician cuando afilamos la sierra y, al mismo tiempo, la sierra se afila cuando tenemos una correcta toma de decisiones y un manejo del tiempo eficaz. Una nos lleva a la otra, mostrando que son parte del mismo fenómeno: una persona integral y altamente efectiva.
	
	\newpage
	
	\begin{thebibliography}{99}
		
		\bibitem{mp1} https://www.gruposura.com/sostenibilidad/
		
		\bibitem{mp2} https://www.epssura.com/
		
		\bibitem{mp3} https://www.gruposura.com/nuestra-compania/etica-y-gobierno-corporativo/
		
		\bibitem{mp4} https://www.gruposura.com/relacion-con-inversionistas/
		
		\bibitem{mp5} https://www.sura.co/somos-sura/nosotros/gobierno-corporativo
		
		\bibitem{mp6} https://www.sura.co/somos-sura/estructura-organizacional
		
		\bibitem{mp7} https://seguros.sura.cl/compania/principios
		
		\bibitem{mp8} https://www.epssura.com/embarazo-menuparalasmamas-229/preguntas-frecuentes-durante-el-embarazo-menuparalasmamas-233?task=section\&id=29
		
		\bibitem{mp9} https://www.epssura.com/valores-corporativos-menu-boletin2-243
		
	\end{thebibliography}
	
	
\end{document}